\documentclass[../principal.tex]{subfiles}
\graphicspath{{\subfix{../imagenes/}}}

\begin{document}

\ifSubfilesClassLoaded{
    \chapter{Electrónica}
}{}

\section{Resistores}

Los resistores son componentes de dos terminales que se utilizan para limitar la corriente eléctrica en un circuito.

Se representan con el símbolo $R$ y se miden en Ohms, cuyo símbolo es la letra griega omega ($\Omega$).

En este libro usamos la palabra resistor para denominar al componente, y la palabra resistencia para referirnos a la propiedad física del material, entonces, el resistor tiene resistencia.

En español muchas veces se usa la palabra resistencia para ambos conceptos, pero preferimos hacer la distinción, tal como se hace en inglés entre resistor y resistance.

\subsection{En serie}

Si tenemos dos resistores en serie, llamados $R_1$ y $R_2$, entre los nodos A y B, podemos diagramarlos así:

\begin{circuitikz}
\draw
  (0,0) node[left]{A}
  to[short, o-] (1,0)
  to[R, l=$R_1$] (3,0)
  to[R, l=$R_2$] (5,0)
  to[short, -o] (6,0) node[right]{B};
\end{circuitikz}

Podemos reemplazar estos dos resistores por uno equivalente $R_{eq}$, que se calcula usando la fórmula:

\begin{equation}
R_{eq} = R_1 + R_2
\end{equation}

y diagramado así:

\begin{circuitikz}
\draw
  (0,0) node[left]{A}
  to[short, o-] (1,0)
  to[R, l=$R_eq$] (3,0)
  to[short, -o] (6,0) node[right]{B};
\end{circuitikz}


\subsection{En paralelo}

Dos resistores en paralelo, llamados $R_1$ y $R_2$, pueden reemplazarse por la resistencia equivalente $R_{eq}$, que se calcula usando la fórmula:


\begin{equation}
R_{eq} = \frac{R_1 \cdot R_2}{R_1 + R_2}
\end{equation}

\newpage

\end{document}

